\chapter{基本概念}
\begin{tcolorbox}[
    colback=blue!5,
    title=自动控制理论的发展
]
    \centering
    {\LARGE \bfseries \color{blue} 自动控制理论发展简史}
    \vspace{1cm}

    \begin{itemize}
        \item[\Large $\bullet$] {\Large \bfseries \color{cyan} 经典控制理论 \quad \small (19世纪初 —— 20世纪50年代)}
        \begin{itemize}
            \item 时域法
            \item 复域法 (根轨迹法)
            \item 频域法
        \end{itemize}
        
        \vspace{0.8cm}
        
        \item[\Large $\bullet$] {\Large \bfseries \color{cyan} 现代控制理论 \quad \small (20世纪60年代 —— )}
        \begin{itemize}
            \item[] \begin{tabular}{lll}
                线性系统 & 自适应控制 & 预测控制 \\
                最优控制 & 鲁棒控制 & 滑模控制 \\
                最佳估计 & 容错控制 & 大系统复杂系统 \\
                系统辨识 & 集散控制 & 非线性系统理论
            \end{tabular}
        \end{itemize}
        
        \vspace{0.8cm}
        
        \item[\Large $\bullet$] {\Large \bfseries \color{cyan} 智能控制理论 \quad \small (20世纪70年代 —— )}
        \begin{itemize}
            \item[] \begin{tabular}{ll}
                专家系统 & 遗传算法 \\
                模糊控制 & 多智能体 \\
                神经网络 & 
            \end{tabular}
        \end{itemize}
    \end{itemize}
    \vspace{0.5cm}
\end{tcolorbox}

\begin{tcolorbox}[colback=blue!6, title=自动控制的概念]
在无人直接参与的情况下，利用控制装置，使工作机械、或生产过程$\color{red}\text{（被控对象）}$的某一个物理量$\color{red}\text{（被控量）}$按预定规律$\color{red}\text{（给定量）}$运行
\end{tcolorbox}

\section*{控制系统方框图示例}

这是一个使用 TikZ 绘制的简单反馈控制系统框图。

\begin{tikzpicture}[
    >=stealth, % 设置全局箭头样式为隐形飞机头
    % 样式定义
    block/.style={rectangle, draw, fill=blue!10, text width=5em, text centered, minimum height=2em},
    % 移除了 node contents，改用传统的 inner sep 和 minimum size
    sum/.style={circle, draw, fill=red!10, inner sep=1pt, minimum size=5mm},
    input/.style={coordinate},
    output/.style={coordinate},
    line/.style={draw, ->}
]

    % 1. 定义节点 (内容统一放在最后的大括号内)
    \node [input] (input) {};
    \node [sum, right=of input, xshift=0.5cm] (sum) {$\times$}; % 内容写在这里
    \node [block, right=of sum, xshift=0.5cm] (controller) {控制器};
    \node [block, right=of controller, xshift=0.5cm] (plant) {受控对象};
    \node [output, right=of plant, xshift=0.5cm] (output) {};
    \node [block, below=of plant, yshift=-0.5cm] (sensor) {反馈环节};

    % 2. 绘制连线
    \draw [line] (input) -- node[above, pos=0.2] {$R(s)$} (sum);
    \draw [line] (sum) -- node[above] {$E(s)$} (controller);
    \draw [line] (controller) -- node[above] {$U(s)$} (plant);
    
    % 绘制到输出，并在中间取一个点用于反馈引出
    \draw [line] (plant) -- node[above, pos=0.8] {$Y(s)$} (output);
    
    % 3. 绘制反馈回路
    % 技巧：先从 plant 和 output 之间引出一条线
    \path (plant) -- (output) node[pos=0.5] (feedback_point) {};
    \draw [line] (feedback_point.center) |- (sensor);
    \draw [line] (sensor) -| (sum) node[left, pos=0.9] {$-$};

    % 在比较点标注正号
    \node at (sum.west) [xshift=-2mm, yshift=2mm] {\tiny $+$};

\end{tikzpicture}

\begin{tcolorbox}[colback=blue!6, title=负反馈原理$\color{red}\text{构成闭环控制系统的核心}$]
将系统的输出信号引回输入端，与输入信号相比较，利用所得的偏差信号进行控制，达到减小偏差，消除偏差的目的。
\end{tcolorbox}

\begin{tcolorbox}[colback=blue!6, title=闭环（反馈）控制系统的特点]
（1）系统内部存在反馈，信号流动构成闭回路\\
（2）偏差起调节作用。
\end{tcolorbox}